\documentclass[12pt, twoside]{article}
%\documentclass[12pt, twoside]{article}
\usepackage[letterpaper, margin=1in, headsep=0.2in]{geometry}
\setlength{\headheight}{0.6in}
%\usepackage[english]{babel}
\usepackage[utf8]{inputenc}
\usepackage{microtype}
\usepackage{amsmath}
\usepackage{amssymb}
%\usepackage{amsfonts}
\usepackage[nomessages]{fp} %\FPeval{\var-name}{2*sin(pi/6)}
\usepackage{siunitx} %units in math. eg 20\milli\meter
\usepackage{yhmath} % for arcs, overparenth command
\usepackage{tikz} %graphics
\usetikzlibrary{quotes, angles, arrows, arrows.meta}
\usepackage{pgfplots}
\usepgfplotslibrary{fillbetween}
\pgfplotsset{compat=1.16}
\usepackage{graphicx} %consider setting \graphicspath{{images/}}
\usepackage{parskip} %no paragraph indent
\usepackage{enumitem}
\usepackage{multicol}
\usepackage{venndiagram}

\usepackage{fancyhdr}
\pagestyle{fancy}
\fancyhf{}
\renewcommand{\headrulewidth}{0pt} % disable the underline of the header
\raggedbottom
\hfuzz=2mm %suppresses overfull box warnings

\usepackage{hyperref}

\fancyhead[LE]{\thepage}
\fancyhead[RO]{Markscheme}
\fancyhead[LO]{La Scuola d'Italia / Dr.\ Huson / IB Math 11 \\* 7 October 2026}

\begin{document}
\subsubsection*{1.10 Quiz: Arithmetic and Geometric Sequences --- Markscheme}

\emph{Calculator allowed. Total: 34 marks.}

\begin{enumerate}[itemsep=0.3cm]

%--- Problem 1: Arithmetic sequence; d, u_n, u_30, show 303 is a term [8 marks] ---
\item[1.] [Maximum mark: 8]

$15,\; 23,\; 31,\; 39,\; \ldots$

\medskip

\textbf{(a)} $d = 23 - 15 = 8$ \hfill \textbf{A1}

\medskip

\textbf{(b)} $u_{n} = u_{1} + (n - 1)d$ \hfill \textbf{(M1)}

$u_{n} = 15 + 8(n - 1)$ \hfill \textbf{A1}

\emph{Accept $u_{n} = 8n + 7$, or any equivalent expression, for full
marks. A candidate who writes $u_{n} = 15 + 8n$ earns M1A0.}

\emph{FT: award A1 for a correct expression from the candidate's own $d$
in (a).}

\medskip

\textbf{(c)} $u_{30} = 15 + 8(30 - 1)$ \hfill \textbf{(M1)}

$= 15 + 232 = 247$ \hfill \textbf{A1}

\emph{A common error is $15 + 8(30) = 255$: M0A0 (it is
$u_{31}$). FT from the candidate's expression in (b).}

\medskip

\textbf{(d)} $15 + 8(n - 1) = 303$ \hfill \textbf{M1}

$8(n - 1) = 288 \;\Rightarrow\; n - 1 = 36 \;\Rightarrow\; n = 37$
\hfill \textbf{A1}

$n$ is a positive whole number, so $303$ is a term: the
$37^{\text{th}}$ term. \hfill \textbf{R1}

\emph{Accept $8n + 7 = 303$, $8n = 296$, $n = 37$. R1 requires the
conclusion linked to $n$ being a whole number; a bare ``$n = 37$'' earns
M1A1R0. Accept a table search that shows $u_{37} = 303$ for full marks.}

\newpage

%--- Problem 2: Two given terms -> u_1 and d; first negative term; terms in k [9 marks] ---
\item[2.] [Maximum mark: 9]

$u_{5} = 38$, \ $u_{14} = 11$.

\medskip

\textbf{(a)} $u_{1} + 4d = 38$ \quad and \quad $u_{1} + 13d = 11$
\hfill \textbf{M1}

$9d = 11 - 38 = -27 \;\Rightarrow\; d = -3$ \hfill \textbf{A1}

$u_{1} = 38 - 4(-3) = 50$ \hfill \textbf{A1}

\emph{M1 for any valid route: both equations, $d = \dfrac{u_{14} -
u_{5}}{14 - 5}$, or solving the system on the GDC. The two A1 marks are
independent. FT: award the second A1 for $u_{1}$ correct from the
candidate's own $d$.}

\medskip

\textbf{(b)} $50 - 3(n - 1) < 0$ \hfill \textbf{M1}

$3(n - 1) > 50 \;\Rightarrow\; n - 1 > 16.67 \;\Rightarrow\; n > 17.67$,
so $n = 18$ \hfill \textbf{A1}

$u_{18} = 50 - 3(18 - 1) = -1$ \hfill \textbf{A1}

\emph{The question asks for the term, not $n$: the final A1 is for $-1$.
A final answer of ``$n = 18$'' alone earns M1A1A0.}

\emph{Accept an equation in place of the inequality, or a table search
showing $u_{17} = 2$ and $u_{18} = -1$ (markscheme-spec §18). FT from the
candidate's $u_{1}$ and $d$.}

\medskip

\textbf{(c)} Equal differences: \ $(4k - 3) - (2k + 1) = (5k + 1) - (4k
- 3)$ \hfill \textbf{M1}

$2k - 4 = k + 4 \;\Rightarrow\; k = 8$ \hfill \textbf{A1}

First three terms: $17,\; 29,\; 41$ \quad ($d = 12$) \hfill \textbf{A1}

\emph{Accept the equivalent set-up $2(4k - 3) = (2k + 1) + (5k + 1)$
(the middle term is the mean of its neighbours). FT: award the final A1
for three terms correctly evaluated from the candidate's $k$.}

\newpage

%--- Problem 3: Arithmetic, geometric or neither; geometric r and u_12; challenge: two given terms [10 marks] ---
\item[3.] [Maximum mark: 10]

\textbf{(a)}

\begin{tabular}{rlll}
(i) & $5,\; 10,\; 20,\; 40$ & geometric, $r = 2$ & \textbf{A1} \\
(ii) & $7,\; 4,\; 1,\; -2$ & arithmetic, $d = -3$ & \textbf{A1} \\
(iii) & $2,\; 3,\; 5,\; 8$ & neither & \textbf{A1} \\
(iv) & $162,\; 54,\; 18,\; 6$ & geometric, $r = \frac{1}{3}$ & \textbf{A1} \\
\end{tabular}

\emph{Each A1 needs the type and, for (i), (ii) and (iv), the correct $d$
or $r$. Accept $r = 0.333$ in (iv). In (iii) the differences are $1, 2,
3$ and the ratios $1.5, 1.67, 1.6$, so neither is constant. $d = 3$ in
(ii) (sign lost) earns A0.}

\medskip

\textbf{(b)} $r = \dfrac{20}{16} = 1.25$ \hfill \textbf{A1}

\medskip

\textbf{(c)} $u_{12} = 16(1.25)^{12 - 1}$ \hfill \textbf{(M1)}

$= 16(1.25)^{11} = 186.26\ldots \approx 186$ \hfill \textbf{A1}

\emph{Accept $186.3$ or $186.26$. $16(1.25)^{12} = 233$ is $u_{13}$:
M0A0. FT from the candidate's $r$ in (b).}

\medskip

\textbf{(d)} From $u_{3}$ to $u_{6}$ is $3$ steps of $\times r$:
$r^{3} = \dfrac{1215}{45} = 27$ \hfill \textbf{M1}

$r = 3$ \hfill \textbf{A1}

$u_{1} = \dfrac{u_{3}}{r^{2}} = \dfrac{45}{9} = 5$ \hfill \textbf{A1}

\emph{Accept the equations $u_{1} r^{2} = 45$ and $u_{1} r^{5} = 1215$,
divided to give $r^{3} = 27$, for the M1. A candidate who uses
$u_{1} = \dfrac{u_{3}}{r} = 15$ earns M1A1A0. FT: award the final A1 for $u_{1}$ correct from the candidate's
own $r$.}

\newpage

%--- Problem 4: Arithmetic vs geometric salary; expressions, year 3, first year B > A [7 marks] ---
\item[4.] [Maximum mark: 7]

Job A: arithmetic, $u_{1} = 38\,000$, $d = 2000$. \ Job B: geometric,
$u_{1} = 36\,000$, $r = 1.06$.

\medskip

\textbf{(a)} Job A: $A_{n} = 38\,000 + 2000(n - 1)$ \hfill \textbf{A1}

Job B: $B_{n} = 36\,000(1.06)^{n - 1}$ \hfill \textbf{A1}

\emph{Accept $A_{n} = 36\,000 + 2000n$. A common error is $r = 0.06$ or
$r = 6$ in Job B: A0. $B_{n} = 36\,000(1.06)^{n}$ earns A0 (it gives
\$$38\,160$ in year $1$).}

\medskip

\textbf{(b)} Job A: $A_{3} = 38\,000 + 2000(3 - 1) = $ \$$42\,000$
\hfill \textbf{A1}

Job B: $B_{3} = 36\,000(1.06)^{3 - 1} = $ \$$40\,449.60$
\hfill \textbf{A1}

\emph{Accept \$$40\,400$ (3 s.f.) or \$$40\,450$. FT: award each A1 for the
value from the candidate's own expression in (a).}

\medskip

\textbf{(c)} Compare the two salaries year by year: \hfill \textbf{(M1)}

\begin{tabular}{ccc}
Year & Job A & Job B \\
\hline
$5$ & \$$46\,000$ & \$$45\,449.17$ \\
$6$ & \$$48\,000$ & \$$48\,176.12$ \\
\end{tabular}
\hfill \textbf{A1}

Year $6$ \hfill \textbf{A1}

\emph{M1 for a valid comparison: a GDC table of both expressions, a
graph, or solving $36\,000(1.06)^{n - 1} = 38\,000 + 2000(n - 1)$ on the
GDC ($n = 5.78\ldots$). The middle A1 is for the two adjacent years (or
the value $n = 5.78$). A single row ($B_{6} > A_{6}$, so year $6$) still
earns full marks, with the feedback note that year $5$ is what justifies
``first'' (markscheme-spec §18).}

\emph{The answer is a whole year: ``year $5.78$'' earns M1A1A0.}

\end{enumerate}
\end{document}
